\bookchapter{Kapitel 30: Bildschirmuebertragung, Praesentation und Auto-Konnektivitaet}

\chapterintro{Auch das größte Handydisplay hat nur 7 Zoll, aber Ihr Fernseher hat 55 und Ihr Beamer 100 Zoll. Das Handybild auf einen großen Bildschirm zu “werfen”, ist eine ganz andere Dimension. Dieses Kapitel erklärt die drei häufigen Szenarien Bildschirmübertragung, Präsentation und Auto-Anbindung und hilft Ihnen, die Fähigkeiten des Handys auf größere Bildschirme auszudehnen.}

\booksection{I. Drei Lager der kabellosen Übertragung —— nicht verwechseln}

Kabellose Übertragung sieht immer so aus: “Handybild auf den Fernseher”. Aber die zugrunde liegenden Technologien sind drei verschiedene, nicht kompatible:

\begin{center}
\begin{tabularx}{\linewidth}{|X|X|X|}
\hline
\textbf{Technologie} & \textbf{Wer nutzt} & \textbf{Merkmale} \tabularnewline
\hline
AirPlay & iPhone/iPad/Mac & Apple-Ökosystem, flüssigste Erfahrung. Benötigt Apple TV oder AirPlay-fähigen Fernseher \tabularnewline
Miracast & Android-Handys & Wi-Fi Direct Direktverbindung, kein Router nötig. Standardmäßig bei Android "kabellose Übertragung" \tabularnewline
DLNA & Alle Plattformen & Push nur Mediendateien (Video/Musik/Bild), kein Spiegeln. Gut, um Videos auf dem TV abzuspielen \tabularnewline
\hline
\end{tabularx}
\end{center}

\booksubsection{Bedienung auf jeder Plattform:}

\begin{itemize}
  \item \textbf{iPhone}: Kontrollzentrum > Bildschirmspiegelung > Ihren Fernseher/Apple TV wählen. Wenn der TV AirPlay nicht unterstützt, können Sie ein Apple TV-Modul oder einen AirPlay-fähigen TV-Stick kaufen.
  \item \textbf{Android}: Benachrichtigungsleiste nach unten ziehen > Schaltfläche "Kabellose Übertragung"/"Smart Cast"/"Multi-Screen-Interaktion" finden > Fernseher wählen. Handy und TV müssen im selben Wi-Fi sein.
  \item \textbf{Universelle Lösung}: Wenn der TV zu alt für kabellose Übertragung ist, kaufen Sie einen Chromecast oder eine Xiaomi-/Huawei-Box und stecken ihn in den HDMI-Anschluss des TVs.
\end{itemize}

\booksection{II. Kabelgebundene Übertragung —— stabiler, null Verzögerung}

Kabellose Übertragung ist bequem, hat aber Verzögerung. Beim Spielen und Präsentieren ist das störend. Kabelgebundene Übertragung löst dieses Problem:

\begin{itemize}
  \item \textbf{iPhone}: Benötigt ein Lightning/USB-C-auf-HDMI-Kabel (ca. 100–200 Yuan). Nach dem Anschließen wird das Handybild direkt an den HDMI-Anschluss des TVs/Beamers ausgegeben.
  \item \textbf{Android}: Einige Handys unterstützen USB-C-auf-HDMI-Direktverbindung (erfordert DP Alt Mode, bei den meisten Flaggschiffen unterstützt). Ein USB-C-auf-HDMI-Kabel genügt.
  \item \textbf{Konferenzraum-Szenario}: Wenn Sie häufig mit dem Handy präsentieren, sollten Sie immer ein HDMI-Kabel dabei haben. Konferenzraum-TVs und Beamer haben fast immer HDMI – anschließen und los, viel zuverlässiger als kabelloses Casting.
\end{itemize}

\booksection{III. Handy wird zum Computer —— Desktop-Modus}

Einige Android-Handys unterstützen den “Desktop-Modus” – nach dem Verbinden mit einem Monitor gibt das Handy eine computerähnliche Desktop-Oberfläche aus, die mit Maus und Tastatur bedient werden kann:

\begin{itemize}
  \item \textbf{Samsung DeX}: Das Aushängeschild der Samsung-Flaggschiffe. Nach dem Verbinden mit einem Monitor erscheint eine Desktop-Oberfläche mit Multi-Window – am nächsten an einem echten Computer. Mit Bluetooth-Tastatur und -Maus können Sie auf Reisen das Notebook zu Hause lassen.
  \item \textbf{Huawei Desktop-Modus}: Huawei Mate-Serie/Pura-Serie bietet nach dem Verbinden mit einem Monitor eine ähnliche Funktion. Unterstützt Multi-Window, Dateiverwaltung und Office-Apps.
  \item \textbf{Motorola Ready For}: Moto-Flaggschiffe unterstützen es ebenfalls, ähnlich wie DeX.
  \item \textbf{sRGB}: Leichtes Büro (E-Mails, Dokumente ändern), PPT präsentieren, Videos auf großem Bildschirm sehen. Nicht geeignet für komplexes Excel oder professionelles Design.
\end{itemize}

\booksection{IV. Auto-Anbindung —— die richtige Art, das Handy im Auto zu nutzen}

Das Bedienen des Handys während der Fahrt ist verboten, aber Navigation, Musik und Telefon sind unverzichtbar. Auto-Anbindungssysteme lassen Sie die Kernfunktionen des Handys auf dem Bordbildschirm nutzen, während das Handy selbst unangetastet bleibt:

\begin{itemize}
  \item \textbf{Apple CarPlay}: Erste Wahl für iPhone-Nutzer. Per USB oder kabellos mit dem Auto verbunden, erscheint auf dem Bordbildschirm eine vereinfachte iPhone-Oberfläche: Navigation (Amap/Baidu Maps), Musik, Telefon, Nachrichten, Podcasts. Siri-Sprachsteuerung vollständig freihändig.
  \item \textbf{Android Auto}: Die Android-Entsprechung. Ähnliche Funktionen wie CarPlay, unterstützt Google Maps, Spotify usw. In China sind aufgrund von Google-Dienst-Beschränkungen einige Funktionen nicht verfügbar.
  \item \textbf{Huawei HiCar}: Huaweis Eigenentwicklung, im chinesischen Markt hohe Unterstützung. Neben Navigation und Musik auch Huawei-Videotelefonie, fusionierte Rückfahrkamera-Bilder und weitere Spezialfunktionen.
  \item \textbf{Baidu CarLife}: Plattformübergreifende Lösung, für iOS und Android. Weniger Funktionen, aber breite Kompatibilität. Viele ältere Modelle unterstützen nur CarLife.
\end{itemize}

\begin{tcbwarn}{Sicherheitshinweis}
Unabhängig vom System sollten Sie beim Fahren immer Sprachsteuerung verwenden. Navigation vor der Fahrt einrichten, Playlist vorher wählen, den Bordbildschirm nicht während der Fahrt bedienen.

\end{tcbwarn}

\booksection{V. Bildschirmübertragung im Konferenzraum —— nicht mehr gestresst}

Besprechungen sind der häufigste Anwendungsfall für Übertragung und auch der fehleranfälligste. Hier sind praktische Tipps:

\begin{itemize}
  \item \textbf{Früh ankommen und testen}: Niemals erst zum Meetingbeginn verbinden. 10 Minuten früh da sein und die Verbindung testen. Wi-Fi-Casting schlägt viel öfter fehl als man denkt.
  \item \textbf{Kabelgebundene Lösung als Plan B}: Ein HDMI-Kabel in der Tasche haben. Wenn kabelloses Casting nicht klappt, ist das Kabel die Rettung.
  \item \textbf{Benachrichtigungsvorschau deaktivieren}: Beim Casting erscheinen Ihre WeChat-Nachrichten auf dem großen Bildschirm. Vor dem Übertragen in den Einstellungen "Benachrichtigungsinhalt auf Sperrbildschirm" deaktivieren oder den "Präsentationsmodus" aktivieren (nicht alle Handys).
  \item \textbf{"Dokument präsentieren" statt "Handy bedienen"}: Wenn Sie ein PPT zeigen, öffnen Sie es mit WPS im Vollbildmodus, statt den ganzen Handy-Desktop zu spiegeln. Das wirkt professioneller und flüssiger.
\end{itemize}

\begin{tcbpractice}{Praxis: Einmal vom Handy auf den Fernseher übertragen}
\begin{enumerate}
  \item Du solltest also vor allem die erste und dritte Kategorie schützen.
  \item Ein Foto oder Video auf dem Handy öffnen.
  \item Kabellose Übertragung versuchen (iPhone: Bildschirmspiegelung, Android: Kabellose Übertragung-Schaltfläche).
  \item Bei Fehlschlag prüfen, ob der TV das Protokoll unterstützt, oder DLNA-Push versuchen (in der Album-/Video-App auf "Cast" tippen).
  \item Nach Erfolg auf dem Handy andere Aktionen ausführen und prüfen, ob der große Bildschirm synchron anzeigt.
\end{enumerate}

\end{tcbpractice}

\begin{tcbsummary}{Kapitelzusammenfassung}
\begin{itemize}
  \item Drei Lager der kabellosen Übertragung: AirPlay (Apple), Miracast (Android), DLNA (plattformübergreifend)
  \item Kabelgebundene Übertragung stabiler und verzögerungsfrei – HDMI-Kabel im Konferenzraum unverzichtbar
  \item Samsung DeX / Huawei Desktop-Modus machen das Handy zum leichten Bürocomputer
  \item Auto-Anbindung: CarPlay / Android Auto / HiCar / CarLife – vier große Lösungen
  \item Eiserne Regel fürs Konferenzraum-Casting: Früh testen, Plan B bereit, Benachrichtigungsvorschau aus
\end{itemize}

\end{tcbsummary}
